\documentclass[border=4pt]{standalone}
\usepackage[siunitx]{circuitikz}

\begin{document}
\begin{circuitikz}[european resistors]
  % Remote, voltage-excited Wheatstone bridge with explicit return.
  \draw[dashed] (-0.55,-0.55) rectangle (3.90,7.15);
  \node[font=\small\bfseries] at (1.68,7.60) {remote bridge sensor};
  \draw (-2.10,0) node[ground] {}
    to[V, l_={$V_{\mathrm{EX}}$}] (-2.10,6.50)
    -- (3.15,6.50);
  \draw (-2.10,0) -- (3.15,0);
  \draw (0.00,6.50)
    to[R, l_={$R_1$}] (0.00,4.65)
    to[R, l_={$R_2$}] (0.00,2.80)
    -- (0.00,0);
  \draw (3.15,6.50)
    to[R, l_={$R_3$}] (3.15,4.65)
    to[R, l_={$R_4$}] (3.15,2.80)
    -- (3.15,0);

  % The two midpoint-signal leads leave the remote boundary as a balanced
  % differential pair; the excitation and return conductors are separate.
  \draw (0.00,4.65) -- (2.75,4.65)
    to[crossing] (3.55,4.65)
    -- (4.65,4.65)
    to[R, l={$R_{F+}=R_F$}] (7.25,4.65)
    -- (11.00,4.65);
  \draw (3.15,2.80) -- (4.65,2.80)
    to[R, l_={$R_{F-}=R_F$}] (7.25,2.80)
    -- (11.00,2.80);
  \node[font=\scriptsize, align=center] at (5.95,6.05)
    {matched series\\input resistors};
  \node[font=\scriptsize, align=center] at (4.90,1.55)
    {balanced\\differential\\signal pair};

  % Differential capacitor and matched common-mode capacitors.
  \draw (8.25,4.65)
    to[C, l_={$C_D$}] (8.25,2.80);
  \draw (9.20,4.65)
    to[C, l={$C_{CM+}$}] (9.20,0);
  \draw (10.15,2.80)
    to[C, l={$C_{CM-}$}] (10.15,0);
  \draw (3.15,0) -- (10.15,0);
  \node[font=\small, align=center] at (8.85,-0.65)
    {$C_{CM+}=C_{CM-}$; filter return};

  % A triangular amplifier symbol shows the functional differential-amplifier
  % behavior. The dashed package boundary carries the INA-specific RG and REF
  % pins that an ordinary op amp does not have.
  \draw[dashed] (11.55,1.10) rectangle (15.25,6.35);
  \node[op amp, noinv input up, scale=1.89, anchor=center]
    (A1) at (13.25,3.725) {};
  \node[font=\scriptsize\bfseries, align=center] at (14.10,5.80)
    {A1\\INA};
  \draw (11.00,4.65) -- (A1.+);
  \draw (11.00,2.80) -- (A1.-);

  % Gain-set pins and resistor terminate at the package boundary.
  \coordinate (RGP) at (12.55,6.35);
  \coordinate (RGM) at (13.45,6.35);
  \draw (RGP) -- (12.55,6.85)
    to[R, l={$R_G$}] (13.45,6.85)
    -- (RGM);
  \node[font=\scriptsize, anchor=north] at (RGP) {RG};
  \node[font=\scriptsize, anchor=north] at (RGM) {RG};

  % Explicit supplies, output load, and output test point.
  \draw (A1.up) -- (14.90,0 |- A1.up)
    -- (14.90,7.20) node[vcc] {$V_{S+}$};
  \draw (A1.down) -- (A1.down |- 0,0.45) node[vee] {$V_{S-}$};
  \draw (A1.out) -- (16.60,3.73)
    node[ocirc, label={[font=\small]right:{TPO: $v_o$}}] {};
  \draw (15.75,3.73)
    to[R, l^={$R_L$}] (15.75,0)
    node[ground, label={[font=\small]right:{TP0: $0$ V}}] {};

  % The REF pin is driven from an explicitly low-impedance source.
  \coordinate (AREF) at (14.25,1.10);
  \draw (AREF) -- (AREF |- 0,1.00)
    -- (14.45,1.00)
    to[V, l={$V_{\mathrm{REF}}$}] (14.45,0)
    node[ground] {};
  \node[font=\scriptsize, anchor=south] at (AREF) {REF};
  \node[font=\small, align=center] at (14.45,-0.70)
    {low-impedance reference};

  % Test points and required junction dots.
  \draw (7.25,4.65) -- (7.25,5.25)
    node[ocirc, label={[font=\small]above:{TP$+$}}] {};
  \draw (7.25,2.80) -- (7.25,2.20)
    node[ocirc, label={[font=\small]below:{TP$-$}}] {};
  \node[circ] at (0.00,6.50) {};
  \node[circ] at (0.00,4.65) {};
  \node[circ] at (0.00,0) {};
  \node[circ] at (3.15,2.80) {};
  \node[circ] at (3.15,0) {};
  \node[circ] at (7.25,4.65) {};
  \node[circ] at (7.25,2.80) {};
  \node[circ] at (8.25,4.65) {};
  \node[circ] at (8.25,2.80) {};
  \node[circ] at (9.20,4.65) {};
  \node[circ] at (10.15,2.80) {};
  \node[circ] at (15.75,3.73) {};
\end{circuitikz}
\end{document}
